\section{Exact adversarial fixtures}\label{sec:fixtures}

Take
\[
 Q=101,\qquad H=8830,\qquad U=99,\qquad h=80.
\]
Exact integer comparisons certify
\[
 8830^{32}\le101^{63}<8831^{32},\qquad
 99^{400}\le101^{399}<100^{400}.
\]
Thus $H=\lfloor Q^{63/32}\rfloor$, $U=\lfloor Q^{399/400}\rfloor$, and $h\le U$:
this is a literal finite floor model of the V59 exponent relations.

For the three shell primes $q=113,127,193$, the cutoff is one.  Their inverses modulo
$80$ are $17,63,17$, respectively, so every row has support
\[
 S_{80,113}=S_{80,127}=S_{80,193}=\{17,63\}.
\]
In particular, residue $63$ is represented by multipliers $-1,1,-1$.  The physical
bucket multiplicity is three.

For uniform amplitudes let the three rows be $v_{113},v_{127},v_{193}$.  Each has
squared norm two, whereas their equal-coefficient sum has squared norm eighteen.
Hence
\begin{equation}
 \frac{\|v_{113}+v_{127}+v_{193}\|^2}
      {\|v_{113}\|^2+\|v_{127}\|^2+\|v_{193}\|^2}=3. \label{eq:ratio-three}
\end{equation}
Equation~\eqref{eq:ratio-three} refutes both physical multiplicity two and a universal
Bessel constant two in this class.  It is a finite structural obstruction, not an
asymptotic lower bound for arithmetic source mass.

A second fixture explains the gcd reduction.  For
$(Q,H,h,a)=(16,65,8,6)$, one has $g=2$ and five rows in bucket $a=6$.  Counting each
multiplier in a class modulo $h$ gives the false upper bound four.  The correct reduced
modulus $h/g=4$ gives upper bound eight and contains all five rows.
